\chapter{Characters and the spectrum they determine}
\label{ch:characters}

A character is a remarkably efficient record of a chiral state space.
It tells us how many states occur at every energy, but it suppresses
their matrix elements and can even fail to distinguish inequivalent
modules. In the present construction, characters become available
after properness, regularity, and the modular data have been
established. We will see how the six-current result selects the
character vector and then turn its coefficients into physical primary
counts.

\section{Three different dimensions attached to a module}

Let $M_i$ be a simple ordinary module, with lowest weight $h_i$.
Write
\begin{equation}
 \chi_i(q)=q^{h_i-1/3}A_i(q),\qquad
 A_i(q)=\sum_{n\ge0}a_i(n)q^n.
 \label{eq:char-expansion}
\end{equation}
The shift $1/3$ is $c/24$ at $c=8$.
The top dimension is $a_i(0)$, an ordinary positive integer.
The categorical quantum dimension is
$d_i=S_{0i}/S_{00}$, generally an irrational algebraic number.
The number of Virasoro primary components at grade $n$ is a third
quantity, obtained after removing descendants. None of these three
numbers can be substituted for another.

For example, the simple module $b$ has top dimension $80$ but
quantum dimension $d^2$, where $d=(3+\sqrt{13})/2$.
The vacuum module has top dimension one, six states at weight one,
and $120$ states at weight two. Only $113$ of the last states are
Virasoro primaries. The remaining seven are the stress tensor and
the derivatives of the six currents.

\section{A hand calculation of the ambient lattice characters}
\index{theta series}\index{character!lattice}

The rank-eight oscillator factor begins
\begin{equation}
 \prod_{m\ge1}(1-q^m)^{-8}
 =1+8q+44q^2+192q^3+\cdots.
 \label{eq:char-eight-oscillators}
\end{equation}
There are eight weight-one oscillators. At weight two there are
eight single oscillators of frequency two and
$\binom{9}{2}=36$ symmetric pairs of frequency-one oscillators.
At weight three the possibilities are one frequency-three oscillator,
one frequency-two and one frequency-one oscillator, or three
frequency-one oscillators. Their counts are
$8+8\cdot8+\binom{10}{3}=192$.

The first theta coefficients of the root lattices are
\[
 \Theta_{E_6}=1+72q+270q^2+720q^3+\cdots,
 \qquad
 \Theta_{A_2}=1+6q+6q^3+\cdots.
\]
Their product is
$1+78q+702q^2+2346q^3+\cdots$. Multiplying by
\eqref{eq:char-eight-oscillators} gives
\begin{equation}
 \chi_B=q^{-1/3}(1+86q+1370q^2+11586q^3+\cdots).
 \label{eq:char-B}
\end{equation}
For instance, the weight-two coefficient is
$702+78\cdot8+44=1370$. Momentum states and oscillators both
contribute, which is why counting roots alone stops working after
weight one.

For $E_8$, the theta series is the weight-four Eisenstein series
\[
 E_4(q)=1+240\sum_{n\ge1}\sigma_3(n)q^n
       =1+240q+2160q^2+6720q^3+\cdots,
\]
where $\sigma_3(n)=\sum_{m\mid n}m^3$.
This follows from the modularity of an even unimodular lattice theta
series and the one-dimensionality of weight-four modular forms for
$SL_2(\Z)$, together with the constant term one. Hence
\begin{equation}
 \chi_{E_8}=\frac{E_4}{\eta^8}
 =q^{-1/3}(1+248q+4124q^2+34752q^3+\cdots).
 \label{eq:char-E8}
\end{equation}
These short computations already supply useful checks on branching.
They do not identify the unknown subalgebra by themselves.

\begin{example}[Why the ambient theory has 248 currents]
\label{ex:char-E8-root-count}
\index{E8 lattice@$E_8$ lattice!root count}
The first character coefficient can also be obtained without a modular
form. In an orthonormal basis $e_1,\ldots,e_8$ of $\R^8$, use
the standard lattice model
\[
 E_8=D_8\cup(D_8+s),\qquad
 D_8=\{n\in\Z^8:\textstyle\sum_i n_i\in2\Z\},
 \qquad s=\tfrac12(e_1+\cdots+e_8).
\]
A current comes from either one of the eight oscillator states or a
momentum of squared norm two. The integer roots are
$\pm e_i\pm e_j$ with $i<j$, giving $4\binom82=112$.
A half-integer vector has squared norm at least $8/4=2$.
Equality requires every coordinate to be $\pm1/2$, and membership
in $D_8+s$ requires an even number of minus signs. Half of the
$2^8$ sign choices have that parity, giving $128$ more roots.
Thus there are $240$ roots and $8+240=248$ currents.

The same model explains the $E_6\oplus A_2$ count. Choose the
$A_2$ simple roots $\beta=e_7+e_8$ and $\gamma=e_6-e_7$;
their Gram matrix has diagonal two and off-diagonal minus one.
A root orthogonal to both obeys
$x_6=x_7=-x_8$. Among integer roots, this forces all nonzero
coordinates into the first five positions, giving
$4\binom52=40$. Among half-integer roots the last three signs
are either $(+,+,-)$ or $(-,-,+)$. In each case exactly $16$
choices of the first five signs give even total parity. Hence the
orthogonal subsystem has $40+32=72$ roots.

To identify its type explicitly, take
\[
 \begin{gathered}
 e_1-e_2,\quad e_2-e_3,\quad e_3-e_4,\quad
 e_4-e_5,\quad e_4+e_5,\\
 \tfrac12(-e_1-e_2-e_3-e_4-e_5+e_6+e_7-e_8).
 \end{gathered}
\]
Their Gram matrix has diagonal entries two and the $E_6$ diagram: the third
root meets the second, fourth, and fifth; the second meets the
first, and the fifth meets the sixth. All other distinct pairings
vanish. The six $A_2$ roots and the $72$ perpendicular roots
therefore contribute $78$ momentum currents to $B$. Adding its
eight oscillators gives $86=(72+6)+(6+2)$, the dimensions of
$\mathfrak e_6\oplus\mathfrak{sl}_3$.

There remain $240-72-6=162$ roots outside this root subsystem.
The product lattice has determinant $3\cdot3=9$, hence index
$\sqrt9=3$ in unimodular $E_8$. The remaining roots lie in the
two nontrivial glue cosets. Negation exchanges those cosets, so each contributes
$81$ currents. This recovers the two glue-top dimensions used
in the branching argument directly from lattice geometry. The
orthonormal coordinates here are for counting; the explicit
generators in Chapter~\ref{ch:lattice_embedding} use their own
simple-root coordinate basis.
\end{example}

\begin{example}[Why a glue top is $27\otimes3$]
\label{ex:char-E8-glue-weights}
The preceding count gives $81$ currents in each glue class.
We can also see the finite representation carried by those currents,
without a table of exceptional Lie-algebra representations.
Keep $\beta=e_7+e_8$ and $\gamma=e_6-e_7$, and choose the
$A_2$ fundamental weight
\[
 \nu=\frac{2\beta+\gamma}{3}
     =\frac{e_6+e_7+2e_8}{3},
 \qquad (\nu,\beta)=1,\quad(\nu,\gamma)=0,\quad \nu^2=\frac23.
\]
An $E_8$ root $r$ projects orthogonally to $\nu$ exactly when
\[
 r_6=r_7,\qquad r_7+r_8=1.
\]
Among the integer roots, the solutions are
$e_8\pm e_i$ for $1\le i\le5$, together with $e_6+e_7$:
eleven roots. Among the half-integer roots, the last three
coordinates must all be $+1/2$. An even number of the first five
signs must be negative, giving $16$ choices. There are therefore
$11+16=27$ roots with this fixed $A_2$ projection.

Their perpendicular components $w=r-\nu$ are $E_6$ weights, all
of squared length
\[
 w^2=r^2-\nu^2=2-\frac23=\frac43.
\]
Each has multiplicity one, since a lattice momentum supplies one
state. For $r=e_6+e_7$, the weight
$w_0=\tfrac23(e_6+e_7-e_8)$ pairs to one with the sixth $E_6$
simple root in the preceding example and to zero with the other
five. It is a fundamental highest weight.

Here is a direct check that these $27$ weights belong to one
representation. Write
\[
 \eta=\tfrac12(-e_1-e_2-e_3-e_4-e_5+e_6+e_7-e_8).
\]
Reflection in $\eta$ sends $w_0$ to the weight with first five
coordinates $+1/2$ and last three $(1/6,1/6,-1/6)$.
The roots $\pm e_i\pm e_j$ in the first five coordinates generate
permutations and even sign changes, producing all sixteen
half-integer choices. Reflecting one with four minus signs in
$\eta$ reaches one of the ten remaining weights; the same
permutations and sign changes reach the other nine. Thus all
$27$ form one Weyl orbit. The root zero modes connect their
one-dimensional weight spaces. This is the irreducible
\emph{minuscule} $27$: its weights form a single Weyl orbit,
each with multiplicity one.

The $A_2$ fundamental representation has the three weights
\[
 \nu,\qquad \nu-\beta,\qquad \nu-\beta-\gamma.
\]
Its Weyl transformations commute with those of $E_6$, so each
of these three projections has the same $27$ perpendicular
weights. They give $27\cdot3=81$ roots in one glue class.
Negation gives the dual representations in the other class.
Including the two adjoint current spaces, the finite branching is
\[
 \mathfrak e_8\big|_{\mathfrak e_6\oplus\mathfrak{sl}_3}
 \cong(78,1)\oplus(1,8)\oplus(27,3)
                         \oplus(\overline{27},\overline3),
 \qquad 248=78+8+81+81.
\]
The choice of which $A_2$ fundamental is called $3$ fixes the
conjugation labels. The two factors of a glue momentum have
conformal weights $2/3$ and $1/3$, because lattice weight is half
the squared length. Their sum is one, explaining why these are
current tops without oscillator excitations.

This is branching under the finite current algebra
$\mathfrak e_6\oplus\mathfrak{sl}_3$. Restriction to the smaller
algebra $U$ still needs the irreducibility argument of
Chapter~\ref{ch:branching}. Neither this count nor equality of
characters specifies the selected Haagerup extension algebra
or its junctions.
\end{example}

\section{Why the principal part can determine the whole vector}
\index{principal part}\index{vector-valued modular function}

For a rational VOA, the characters form a vector-valued modular
function. The $T$ matrix contains the phase
$e^{2\pi i(h_i-c/24)}$, not merely the categorical twist
$e^{2\pi ih_i}$. The difference matters when deciding whether a
constant vector is modular invariant.

The \emph{principal part} is the sum of all negative powers of
$q$ in the character vector. In the Haagerup modular representation,
strictly positive nonvacuum lowest weights leave only two possible
polar terms:
\[
 q^{-1/3}\quad\text{in the vacuum coordinate},\qquad
 \gamma q^{-7/39}\quad\text{in the $d_3$ coordinate}.
\]
The second exponent would arise from a lowest weight $2/13$.
The integer $\gamma\ge0$ is its top dimension if that weight occurs;
$\gamma=0$ means that this sector begins one or more grades later.

\begin{proposition}[Uniqueness for the specified principal part]
\label{prop:char-unique}
For the Haagerup character representation at $c=8$, two character
vectors with the same $\gamma$ agree in every coefficient.
\end{proposition}
\begin{proof}
The modular representation has finite image, a consequence of the
congruence theorem for the rational VOA modular data used here.
Its kernel $\Gamma$ therefore has finite index in $SL_2(\Z)$.
Subtract the two vectors. Every component of the difference is
holomorphic at infinity and, by its modular transformation law,
at all cusps of $\Gamma$. A cusp is a $\Gamma$-orbit of rational
boundary points or infinity. The quotient
$\Gamma\backslash\mathfrak H$ of the upper half-plane becomes a
compact Riemann surface after adding these finitely many cusps.
Each component is holomorphic on $\mathfrak H$, and the absence of
negative powers in its local $q$-expansion at every cusp means that
it extends holomorphically over the added points. The maximum
principle then makes it constant on this compact modular curve.
The resulting vector must be fixed by both $S$ and $T_{\rm ch}$.

Only the coordinate with categorical twist $e^{2\pi i/3}$ is fixed
by $T_{\rm ch}=e^{-2\pi i/3}\Theta$. But its $S$-image has a
nonzero vacuum component, so no nonzero vector supported on that
coordinate is also fixed by $S$. The constant difference is zero.
\end{proof}

Existence of the relevant vector-valued modular functions is a
separate input. The Evans--Gannon construction supplies fundamental
solutions $\Xi_{:,0}$ and $\Xi_{:,d_3}$ with the two allowed
principal parts; the colon denotes the full twelve-component column.
Their low coefficients give
\begin{equation}
 \boldsymbol\chi=\Xi_{:,0}+\gamma\Xi_{:,d_3},
 \qquad \dim U_1=6+13\gamma.
 \label{eq:char-gamma}
\end{equation}
The independently proved six-current result forces $\gamma=0$.
Together with Proposition~\ref{prop:char-unique}, this selects the
entire character vector. The logic uses much more than matching a
few coefficients: modular covariance, the allowed poles, absence of
invariant constants, and a proved dimension of $U_1$ all enter.

This also clarifies the order of the construction. The first few
vacuum dimensions $1,6,120,956$ are outputs of finite lattice
closure. The subsequent coefficient $6010$ is supplied by the
character theorem. Using that coefficient to prove the earlier
closure statement would reverse the dependency.

\section{The twelve characters at their first grades}
\index{character!Haagerup sectors}

The selected vector has the following initial data. Every row
represents a simple extended-algebra module, and the grade is
measured above its actual lowest weight.
\begin{center}\small
\begin{tabular}{c|c|rrrr}
\toprule
Module&$h_i$&$a_i(0)$&$a_i(1)$&$a_i(2)$&$a_i(3)$\\
\midrule
$0$&$0$&1&6&120&956\\
$b$&$1$&80&1250&10630&65042\\
$a$&$1$&81&1377&11583&71037\\
$c_0$&$1$&81&1377&11583&71037\\
$c_1$&$1/3$&3&243&2916&21870\\
$c_2$&$2/3$&27&594&5967&39852\\
$d_1$&$6/13$&7&292&3204&23010\\
$d_2$&$11/13$&42&777&7147&45367\\
$d_3$&$15/13$&119&1623&12996&76429\\
$d_4$&$5/13$&5&229&2738&19942\\
$d_5$&$7/13$&13&347&3804&26390\\
$d_6$&$8/13$&14&441&4445&30329\\
\bottomrule
\end{tabular}
\end{center}

There are two instructive features. First, $d_3$ has twist exponent
$2/13$ but lowest energy $15/13$. Twists know energies only modulo
one. Second, $a$ and $c_0$ are different simple modules with identical
unflavoured characters. A trace of $q^{L_0-c/24}$ forgets enough
structure to make this possible. Their defect eigenvalues and their
categorical labels still distinguish them.

\begin{example}[Branching checked grade by grade]
The restriction $B\downarrow_U=U\oplus b$ implies
$\chi_B=\chi_0+\chi_b$. Using the table gives
\[
 86=6+80,\qquad1370=120+1250,\qquad
 11586=956+10630,
\]
in agreement with the independent lattice calculation.
For the selected holomorphic extension,
\begin{equation}
 \chi_{E_8}=\chi_0+\chi_b+2\chi_a.
 \label{eq:char-branch-E8}
\end{equation}
The first three nontrivial checks are
\[
 \begin{split}
 248&=6+80+2\cdot81,\\
 4124&=120+1250+2\cdot1377,\\
 34752&=956+10630+2\cdot11583.
 \end{split}
\]
Replacing $a$ by $c_0$ leaves the scalar character identity unchanged.
Thus the identity verifies graded dimensions once branching is known;
it cannot select an extension algebra or prove its multiplication.
\end{example}

\section{Removing Virasoro descendants}
\index{primary multiplicity}\index{descendant subtraction}

Let $p(n)$ be the integer partition number. A positive-weight
irreducible Virasoro Verma module has descendant generating function
$\prod_{m\ge1}(1-q^m)^{-1}$. At $c=8$, every positive real highest
weight has such an irreducible Verma module. Here is a useful check
of this otherwise nontrivial statement.

Parametrize $c=13-6(t+t^{-1})$. For $c=8$, $t$ lies on the unit
circle with $t+t^{-1}=5/6$ and nonzero imaginary part. The Kac
determinant is the general-level analogue of the Virasoro Gram
determinant computed in Chapter~\ref{ch:cft}. Its possible singular
weights are
\[
 h_{r,s}=\frac14\bigl((r^2-1)t+(s^2-1)t^{-1}-2rs+2\bigr),
 \qquad r,s\in\Z_{>0}.
\]
Reality requires $r=s$. Then
$h_{r,r}=-7(r^2-1)/24\le0$. No positive real highest weight is
singular. At weight zero, the primitive vacuum null is
$L_{-1}\one$, which has already been removed in the vacuum module.

To turn this irreducible-Verma calculation into a count of actual
primary components, use the strong-unitarity conclusion proved in
Chapter~\ref{ch:unitarity}. In each finite-dimensional energy space,
take the orthogonal complement of the descendants of lower-weight
vectors. Repeating in increasing energy decomposes each module into
Virasoro highest-weight summands. For the vacuum module, the needed
positive form is already inherited from the lattice.

Put $P(q)=\prod_{m\ge1}(1-q^m)$. Let $v_i(n)$ count Virasoro
primary components of weight $h_i+n$, including the vacuum once.
For nonvacuum modules,
\begin{equation}
 \sum_{n\ge0}v_i(n)q^n=P(q)A_i(q),\qquad i\ne0.
 \label{eq:char-primary-nonvac}
\end{equation}
For the vacuum module the correct formula is
\begin{equation}
 \sum_{n\ge0}v_0(n)q^n=P(q)A_0(q)+q.
 \label{eq:char-primary-vac}
\end{equation}
To understand the correction, the vacuum descendants contribute
$(1-q)/P(q)$ rather than $1/P(q)$. If $B_0(q)$ counts all vacuum
module primaries, then $A_0=(B_0-q)/P$, hence $B_0=PA_0+q$.

\begin{example}[The first interacting primary space]
Using
$P(q)=1-q-q^2+O(q^5)$ and
$A_0=1+6q+120q^2+956q^3+6010q^4+\cdots$, one obtains
\[
 \sum_nv_0(n)q^n
 =1+6q+113q^2+830q^3+4934q^4+\cdots.
\]
At weight two the subtraction is
$120-6-1=113$: six current derivatives and one stress.
At weight three it is $956-120-6=830$. The generating-function
method automatically accounts for descendants of descendants.
For $c_1$, the first new grade-one primaries number
$243-3=240$, since each of its three top states has one
$L_{-1}$ descendant.
\end{example}

The same calculation for $E_8$ gives $248$ weight-one primaries and
$4124-248-1=3875$ weight-two primaries. These are component counts,
not counts of affine highest-weight multiplets. A whole multiplet
can contribute many Virasoro primaries of the same weight.

\section{Which deductions require more than a character}

A character establishes a graded count once the module is known to
exist. It does not give a preferred orthonormal basis, the action of
the currents, the reduced matrix elements of an OPE, or a defect
junction associator. In particular, the coefficient $113$ predicts
the dimension of a space on which a three-point tensor can be
computed. It does not specify that tensor.

A finite truncation also has a precise coverage limit. If all
characters are known through grade $N$ and all lowest weights are
nonnegative, then the diagonal full spectrum is complete through
total dimension $N$: every contributing grade is at most $N$.
The rectangular product of those truncated series contains some
states at larger dimensions, but is generally incomplete there.
Reporting the larger partial rectangle as a complete spectrum would
be an error even if every retained coefficient were exact.

There is a useful division of labor. The lattice construction gives
actual fields and products. The structural argument gives modules
and modular data. The character theorem gives every unflavoured
graded dimension. Descendant subtraction gives Virasoro-primary
counts. Finally, a choice of full theory specifies how holomorphic
and antiholomorphic pieces combine. We carry out that last step in
Chapter~\ref{ch:full_theories}.

\section{Exercises}

\begin{exercise}\label{ex:char-oscillator}
Derive the coefficient $726$ of $q^4$ in the rank-eight oscillator
factor by classifying the partitions $4$, $3+1$, $2+2$, $2+1+1$,
and $1+1+1+1$. Use it and $E_4$ to obtain the weight-four
coefficient of $V_{E_8}$.

\par\smallskip\noindent{\footnotesize\normalfont
\hyperref[sol:ex:char-oscillator]{Solution on page~\pageref*{sol:ex:char-oscillator}.}}\par
\end{exercise}

\begin{exercise}\label{ex:char-primary}
Compute the first three primary multiplicities above the top of
the $b$ module and of the $d_4$ module. Reconstruct the corresponding
character coefficients by convolving with partition numbers.

\par\smallskip\noindent{\footnotesize\normalfont
\hyperref[sol:ex:char-primary]{Solution on page~\pageref*{sol:ex:char-primary}.}}\par
\end{exercise}

\begin{exercise}\label{ex:char-pole}
Explain why the categorical twist $e^{4\pi i/13}$ alone permits
lowest weights $2/13,15/13,28/13,\ldots$. Identify exactly which
steps in this chapter select $15/13$.

\par\smallskip\noindent{\footnotesize\normalfont
\hyperref[sol:ex:char-pole]{Solution on page~\pageref*{sol:ex:char-pole}.}}\par
\end{exercise}

\begin{exercise}\label{ex:char-equal}
Two distinct simple modules have equal unflavoured characters.
Explain why this does not violate linear independence of simple
objects in the Grothendieck group. Give two types of observables
that could distinguish the modules.

\par\smallskip\noindent{\footnotesize\normalfont
\hyperref[sol:ex:char-equal]{Solution on page~\pageref*{sol:ex:char-equal}.}}\par
\end{exercise}

\begin{exercise}\label{ex:char-uniqueness-scope}
In Proposition~\ref{prop:char-unique}, replace the modular
representation by one containing a trivial summand. Explain why
the same principal part need no longer determine a unique function.
What extra normalization would be required?

\par\smallskip\noindent{\footnotesize\normalfont
\hyperref[sol:ex:char-uniqueness-scope]{Solution on page~\pageref*{sol:ex:char-uniqueness-scope}.}}\par
\end{exercise}
